\section{Discusión}

\subsection{Analogía limitada con sistemas de múltiples escalas}

La dinámica entre capas puede visualizarse, de forma heurística y limitada, mediante la analogía de un huracán maduro. Las intensificaciones locales (celdas convectivas intensas) corresponden a presentes locales intensificados. La organización entre escalas (bandas, eyewall, circulación general) corresponde a la reducción de fragmentación y al aumento de coherencia entre presentes locales. La inercia del huracán maduro ilustra el efecto de filtrado probabilístico de las capas superiores frente a intensificaciones locales fragmentadas. La desorganización (cizalladura, interacción con tierra) corresponde al aumento de antisincronización.

La analogía es útil para transmitir que la intensificación local no equivale automáticamente a cambio estructural, y que la organización entre escalas es condición necesaria para que las intensificaciones locales tengan consecuencias globales. Sin embargo, presenta limitaciones importantes: en el huracán la organización tiene un fuerte componente físico y energético; en el presente, la coherencia es fundamentalmente relacional. Además, el huracán posee un ciclo de vida claro, mientras que no está establecido si el presente de capas superiores sigue un patrón análogo de consolidación, madurez y posible decadencia o fragmentación creciente.

\subsection{Preguntas abiertas}

El marco aquí propuesto deja abiertas varias líneas de investigación:

\begin{enumerate}
    \item ¿En qué condiciones precisas una intensificación del presente local logra superar el efecto de filtrado probabilístico ascendente («resistencia ascendente») de las capas superiores y generar reorganización estructural efectiva?
    
    \item ¿Existe un análogo de ``ciclo de vida'' del presente de capas superiores (consolidación, madurez, posible decadencia o aumento de fragmentación)?
    
    \item ¿Cómo se comporta la dinámica de fragmentación y filtrado probabilístico en sistemas con acoplamiento fuertemente asimétrico (un módulo con peso dominante)?
    
    \item ¿Es posible detectar empíricamente casos de retrocausalidad local, entendida como reversión hacia configuraciones estructurales anteriores mediada por alta antisincronización sostenida?
    
    \item ¿Qué rol preciso juega la memoria activa (trapping + hiper-persistencia) en el efecto de filtrado probabilístico descendente («resistencia descendente») frente a intensificaciones locales fragmentadas?
    
    \item ¿Cómo se articula esta visión estratificada del presente con concepciones filosóficas del tiempo de la esencia, en particular con el proyecto antropológico de Leonardo Polo (véase \citep{polo_antropologia_trascendental,padilla_polo_dialogue_2026})?
\end{enumerate}

Estas preguntas no son meramente técnicas; señalan los puntos en los que la ontología del presente debe ser refinada o, eventualmente, corregida a la luz de nueva evidencia empírica o de mayor profundidad conceptual.

\subsection{Limitaciones y objeciones al marco propuesto}

El marco es exploratorio y reconoce las siguientes limitaciones y objeciones legítimas de forma directa.

\textbf{Riesgo de reificación de las «capas».} Las capas se definen mediante marcadores operacionales. El marco no afirma que sean entidades sustanciales separadas de los procesos de persistencia; se trata de una jerarquía analítica débil. Si la discreción observada depende del umbralado o de la resolución elegida, la partición categorial requerirá revisión.

\textbf{Carácter heurístico del término «resistencia».} El concepto operativo central es el de filtrado probabilístico (efecto estructural sin agente). Los términos «resistencia ascendente/descendente» se mantienen entre comillas únicamente como etiquetas heurísticas secundarias. Su uso puede refinarse o abandonarse en favor de «filtrado probabilístico ascendente/descendente».

\textbf{Ausencia de mecanismo de causación descendente.} El marco describe una reducción de probabilidad de impacto de fluctuaciones locales cuando la Capa 2 está consolidada, pero no especifica el mecanismo causal preciso (retroalimentación, topología del espacio de estados u otro). Esta es una limitación actual que exige modelado adicional o diseños experimentales de intervención.

\textbf{Posible circularidad entre métricas y ontología.} Las definiciones se apoyan en las métricas que detectan los fenómenos. El marco sostiene que las capas introducen propiedades nuevas (coherencia de escalas; reorganización global de la matriz), pero la distinción entre descripción métrica y afirmación ontológica sigue requiriendo elaboración filosófica adicional.

\textbf{Subdeterminación de la ontología por los datos.} Los datos provienen de sistemas sintéticos controlados y casos reales limitados. Otras ontologías podrían acomodar los mismos patrones. El marco se ofrece como una síntesis coherente con las observaciones; su valor dependerá de la generación de predicciones contrastables y del diálogo con tradiciones filosóficas relevantes (como la de Leonardo Polo). La subdeterminación actual es reconocida.